\documentclass[UTF8,a4paper,11pt]{ctexart}

% ── Page geometry ──
\usepackage[margin=2.2cm,top=2.5cm,bottom=2.5cm]{geometry}

% ── Typography ──
\usepackage{fontspec}
\usepackage{setspace}
\setstretch{1.2}

% ── Graphics ──
\usepackage{graphicx}
\graphicspath{{/home/zhiwen/projects/mooc2handout-skill/coursera-mcp/module-4/frames/}}

% ── Colors ──
\usepackage{xcolor}
\definecolor{accent}{HTML}{1a73e8}
\definecolor{darkgray}{HTML}{333333}
\definecolor{lightgray}{HTML}{f5f5f5}
\definecolor{keyframebg}{HTML}{fafafa}
\definecolor{videoblue}{HTML}{065fd4}
\definecolor{notebg}{HTML}{e8f0fe}
\definecolor{warnbg}{HTML}{fef7e0}

% ── Links ──
\usepackage{hyperref}
\hypersetup{
  colorlinks=true,
  linkcolor=accent,
  urlcolor=videoblue,
  citecolor=accent,
  pdftitle={模型上下文协议导论},
  pdfauthor={}
}

% ── Layout ──
\usepackage{enumitem}
\setlist[itemize]{leftmargin=1.8em,itemsep=2pt}
\setlist[enumerate]{leftmargin=1.8em,itemsep=2pt}
\usepackage{booktabs}
\usepackage{longtable}
\usepackage{tcolorbox}
\usepackage{titlesec}
\usepackage{fancyhdr}
\usepackage{lastpage}
\usepackage{amssymb}

% ── Header/Footer ──
\pagestyle{fancy}
\fancyhf{}
\fancyhead[L]{\small\textcolor{darkgray}{模型上下文协议导论}}
\fancyhead[R]{\small\textcolor{darkgray}{第四单元：课程回顾}}
\fancyfoot[C]{\small\textcolor{darkgray}{\thepage\ / \pageref{LastPage}}}
\renewcommand{\headrulewidth}{0.4pt}

% ── Section styling ──
\titleformat{\section}
  {\Large\bfseries\color{accent}}
  {\thesection}{1em}{}
\titleformat{\subsection}
  {\large\bfseries\color{darkgray}}
  {\thesubsection}{1em}{}

% ── Custom environments ──
\newtcolorbox{keyconcept}{
  colback=notebg, colframe=accent,
  fonttitle=\bfseries, boxrule=0.5pt,
  left=8pt, right=8pt, top=6pt, bottom=6pt
}
\newtcolorbox{supplement}{
  colback=lightgray, colframe=darkgray,
  fonttitle=\bfseries, boxrule=0.3pt,
  left=8pt, right=8pt, top=6pt, bottom=6pt
}
\newtcolorbox{exercise}{
  colback=warnbg, colframe=orange,
  fonttitle=\bfseries, boxrule=0.5pt,
  left=8pt, right=8pt, top=6pt, bottom=6pt
}

% ── Video link command ──
\newcommand{\videolink}[2]{%
  \href{#1}{\textcolor{videoblue}{\textbf{#2}}}%
}

% ── Keyframe figure command ──
\newcommand{\keyframe}[2]{%
  \begin{center}
    \fcolorbox{lightgray}{keyframebg}{%
      \includegraphics[width=0.75\linewidth]{#1}%
    }
    \par\vspace{2pt}
    {\small\textcolor{darkgray}{#2}}
  \end{center}
}

\begin{document}

\begin{center}
  \vspace*{2cm}
  {\Huge\bfseries\color{accent} 模型上下文协议导论}\\[0.3cm]
  {\small\textcolor{darkgray}{Introduction to Model Context Protocol}}\\[0.8cm]
  {\LARGE 第四单元：课程回顾}\\[0.3cm]
  {\small\textcolor{darkgray}{Module 4: Review}}\\[1.2cm]
  {\large 授课教师： }\\[0.5cm]
  {\large\textcolor{darkgray}{由 mooc2handout 自动生成}}\\[0.3cm]
  {\small\textcolor{darkgray}{\today}}
  \vfill
  {\small 本讲义由课程字幕、补充材料与可选的 AI 推断关键帧自动生成。}
\end{center}
\newpage

\tableofcontents
\newpage

\section*{本单元知识地图}
\addcontentsline{toc}{section}{本单元知识地图}
\begin{longtable}{p{3.1cm}p{10.2cm}}
\toprule
主题 & 核心问题 \\
\midrule
控制方 & 谁决定 tool / resource / prompt 什么时候被调用？ \\
选型规则 & 什么场景该用哪种原语？ \\
回顾目标 & 整门课最终想建立什么判断框架？ \\
\bottomrule
\end{longtable}

\section*{0. 概要}
\addcontentsline{toc}{section}{0. 概要}
最后一单元把整门课收束到一个简单但实用的判断框架：当能力应该由模型决定时用 tools，
当能力应该由应用决定时用 resources，当能力应该由用户决定时用 prompts。

\begin{itemize}
  \item 先问“谁来触发”，再问“该用哪种原语”。
  \item 这一套分类能直接指导应用设计。
  \item MCP 的价值在于把能力暴露方式讲清楚。
\end{itemize}

\section*{1. MCP 回顾}
\addcontentsline{toc}{section}{1. MCP 回顾}
\noindent\textbf{回看：}\videolink{https://www.coursera.org/learn/introduction-to-model-context-protocol/home/module/4}{MCP 回顾}

\textbf{本讲重点：} 把三类原语收束成一个选择框架。

这一讲重新回顾 tools、resources、prompts 的差别，并强调一个最重要的判断方法：
看看到底是谁在决定它什么时候被调用。

如果是模型在决定，就更像 tools；如果是应用在决定，就更像 resources；
如果是用户在决定，就更像 prompts。这个判断框架比死记术语更有用。

\begin{itemize}
  \item tools 是 model-controlled。
  \item resources 是 app-controlled。
  \item prompts 是 user-controlled。
\end{itemize}

\begin{keyconcept}
\textbf{自动摘要}
根据字幕自动扩写，这一讲再补充几条最值得记住的线索。
\begin{itemize}
  \item \textbf{形式定义}: 这一讲梳理 AND / OR / NOT / IF...THEN 等连接词的真值条件，并提醒不要把日常语言的直觉直接搬进数学里。
\end{itemize}
\end{keyconcept}

\section*{2. 本单元最该记住的五句话}
\addcontentsline{toc}{section}{2. 本单元最该记住的五句话}
\begin{enumerate}
  \item tools 是 model-controlled。
  \item resources 是 app-controlled。
  \item prompts 是 user-controlled。
  \item 选择原语时先看触发权归谁。
  \item MCP 的核心价值是让能力暴露方式清晰、统一、可复用。
\end{enumerate}

\section*{3. 练习题}
\addcontentsline{toc}{section}{3. 练习题}
\subsection*{A. 选择题（5题）}
\begin{enumerate}
  \item tools 属于哪种控制方式？\\
  A. model-controlled \quad B. app-controlled \quad C. user-controlled \quad D. OS-controlled

  \item resources 属于哪种控制方式？\\
  A. model-controlled \quad B. app-controlled \quad C. user-controlled \quad D. 随机控制

  \item prompts 属于哪种控制方式？\\
  A. model-controlled \quad B. app-controlled \quad C. user-controlled \quad D. 网络控制

  \item 如果你想让应用根据用户点击来触发一段固定指令模板，优先用什么？\\
  A. tools \quad B. prompts \quad C. resources \quad D. database

  \item MCP 的一个重要价值是什么？\\
  A. 让每个应用都写一套私有协议 \quad B. 把能力暴露方式结构化 \quad C. 取消 client \quad D. 只支持一种传输

\end{enumerate}

\subsection*{B. 判断题（3题）}
\begin{enumerate}
  \item 选择原语时，先判断由谁来决定它被触发。（\quad）
  \item resources 更适合应用主动注入上下文数据。（\quad）
  \item prompts 的触发权通常在用户手里。（\quad）
\end{enumerate}

\subsection*{C. 简答题（2题）}
\begin{enumerate}
  \item 分别用一句话概括 tools、resources、prompts 的职责。
  \item 举一个你会优先使用 resources 而不是 tools 的场景。
\end{enumerate}

\section*{4. 参考答案}
\addcontentsline{toc}{section}{4. 参考答案}
\subsection*{选择题答案}
1.A \quad 2.B \quad 3.C \quad 4.B \quad 5.B

\subsection*{判断题答案}
1. 对 \quad 2. 对 \quad 3. 对

\subsection*{简答题要点}
\begin{enumerate}
  \item tools 由模型决定何时调用；resources 由应用决定何时提供上下文；prompts 由用户决定何时触发模板。
  \item 例如用户引用一份文档并希望自动把内容注入上下文时，更适合用 resources。
\end{enumerate}

\end{document}